\begin{abstract}
    \textbf{When discussing state-sponsored cyber actors, Russia is often among the first countries associated with this topic. Its cyber operations are run by  an ecosystem of state-employed operators, state-sponsored hacking groups and cybercriminals working with the government’s approval \parencite{olejnik2025russian}. At the core of this ecosystem are three intelligence services. The GRU, Russia’s military intelligence, is known for its aggressive cyber operations. The FSB is officially in charge of fighting cybercrime, but in practice it uses criminal networks to expand its intelligence collection while keeping deniability. The SVR focuses on long-term foreign intelligence collection and persistent espionage. These three services also compete with each other \parencite{safebreach-russia-guide}. Each runs its own cyber units, which Western authorities track under different names: Sandworm is linked to the GRU, Star Blizzard is believed to be subordinate to the FSB, and APT29, also known as Cozy Bear, almost certainly belongs to the SVR \parencite{cisa-russia-advisories}. More broadly, Russian information warfare aims to weaken support for Western policies and to divide NATO \parencite{olejnik2025russian}. Since the three intelligence services have different tasks and methods, it is useful to analyse one of their units in more detail. This report therefore focuses on APT29 as an example of a state-sponsored actor. Its long-term espionage activities show well, how Russian cyber operations combine technical skills with an understanding of human behaviour.}
\end{abstract}

